% =============================================================================
% Release 2 — Sprint 2 : Intégration continue et déploiement continu (CI/CD)
% =============================================================================

\section[Sprint 2 : Intégration continue et déploiement continu (CI/CD)]%
{Sprint 2 : Intégration continue et déploiement continu (CI/CD)}

Ce sprint industrialise la validation du code, la construction des artefacts conteneurisés
et l'automatisation du provisionnement infrastructurel. Il s'appuie sur \textbf{Azure DevOps}
et un pipeline entièrement décrit en YAML (\texttt{azure-pipelines.yml}).


% ---------------------------------------------------------------------------
\subsection{Objectif du Sprint 2}
% ---------------------------------------------------------------------------

Le tableau~\ref{tab:r2-s2-objectifs} résume les objectifs du Sprint~2 par axe.

\begin{table}[H]
\centering
\caption{Objectifs du Sprint~2 par axe}
\label{tab:r2-s2-objectifs}
\footnotesize
\renewcommand{\arraystretch}{1.1}
\setlength{\tabcolsep}{6pt}
\begin{tabular}{|p{2.8cm}|p{9.2cm}|}
\hline
\textbf{Axe} & \textbf{Objectif} \\
\hline
Intégration continue &
  Pipeline déclenché sur \texttt{main} et \texttt{develop}, exécuté sur agent
  auto-hébergé (\texttt{ticketmanager-pool}) \\
\hline
Qualité &
  Analyse statique SonarQube ; tests automatisés backend (Jest) \\
\hline
Sécurité &
  Scan des images Docker (Trivy) ; analyse des dépendances (Snyk, optionnel) \\
\hline
Conteneurisation &
  Images multi-étapes \texttt{ticket-backend} et \texttt{ticket-frontend} \\
\hline
Registre &
  Publication des images vers Nexus Repository (\texttt{localhost:8082}) \\
\hline
IaC automatisé &
  Exécution Terraform (init, plan, apply) intégrée au pipeline \\
\hline
\end{tabular}
\end{table}

\paragraph{Résultat attendu.}
\begin{itemize}
  \item pipeline Azure DevOps vert de bout en bout (stages Build à Push minimum) ;
  \item images Docker versionnées dans Nexus (\texttt{docker-ticketmanager/}) ;
  \item rapport SonarQube pour le projet \texttt{ticket-manager} ;
  \item infrastructure Terraform synchronisée via le stage IaC.
\end{itemize}

\paragraph{Note sur les stages temporairement désactivés.}
Dans la version actuelle de \texttt{azure-pipelines.yml}, les stages \textbf{Test},
\textbf{Snyk} et \textbf{Trivy} sont commentés pour accélérer les itérations. Le squelette
du rapport prévoit leurs figures ; réactivez-les avant la soutenance si possible.


% ---------------------------------------------------------------------------
\subsection{Backlog du Sprint 2}
% ---------------------------------------------------------------------------

Le tableau~\ref{tab:r2-s2-backlog} présente le backlog planifié pour le Sprint~2.

\subsubsection{Epic — Chaîne d'intégration continue}

\begin{table}[H]
\centering
\caption{Backlog Sprint~2 — Epic intégration continue et déploiement continu}
\label{tab:r2-s2-backlog}
\footnotesize
\renewcommand{\arraystretch}{0.85}
\setlength{\tabcolsep}{4pt}
\begin{tabular}{|c|p{7.8cm}|c|c|}
\hline
\textbf{ID} & \textbf{User story} & \textbf{Priorité} & \textbf{Est. (j)} \\
\hline
R2-S2-C01 &
  En tant que développeur, je centralise le code source sur Azure Repos &
  H & 1 \\
\hline
R2-S2-C02 &
  En tant que développeur, j'exécute le pipeline sur un agent auto-hébergé &
  H & 2 \\
\hline
R2-S2-C03 &
  En tant que développeur, je déclenche automatiquement le pipeline à chaque push &
  H & 1 \\
\hline
R2-S2-C04 &
  En tant que développeur, je compile le backend NestJS et le frontend Next.js à chaque livraison &
  H & 2 \\
\hline
R2-S2-C05 &
  En tant que développeur, j'exécute les tests automatisés et publie la couverture &
  H & 2 \\
\hline
R2-S2-C06 &
  En tant que développeur, j'intègre l'analyse SonarQube dans le pipeline &
  H & 2 \\
\hline
R2-S2-C07 &
  En tant que développeur, je construis des images Docker prêtes pour la production &
  H & 2 \\
\hline
R2-S2-C08 &
  En tant que développeur, je scanne les images Docker avec Trivy avant publication &
  M & 1 \\
\hline
R2-S2-C09 &
  En tant que développeur, je publie les images validées dans Nexus Repository &
  H & 1 \\
\hline
R2-S2-C10 &
  En tant que développeur, j'automatise Terraform dans le pipeline avec validation manuelle &
  H & 2 \\
\hline
\end{tabular}
\end{table}


% ---------------------------------------------------------------------------
\subsection{Analyse — Conception du pipeline}
% ---------------------------------------------------------------------------

\subsubsection{Présentation d'Azure DevOps}

Azure DevOps regroupe les outils collaboratifs mobilisés dans TicketManager :
\begin{itemize}
  \item \textbf{Azure Boards} — planification agile et suivi des sprints ;
  \item \textbf{Azure Repos} — gestion de versions Git ;
  \item \textbf{Azure Pipelines} — exécution automatisée du pipeline CI/CD en YAML.
\end{itemize}

La figure~\ref{fig:r2-s2-ado-ui} présente l'interface principale d'Azure DevOps utilisée
dans le projet.

\begin{figure}[htbp]
    \centering
    \includegraphics[width=0.85\textwidth]{s2-azure-devops-ui.png}
    \caption{Interface principale d'Azure DevOps}
    \label{fig:r2-s2-ado-ui}
\end{figure}

\subsubsection{Structure déclarative du pipeline}

Le pipeline est organisé en \textbf{stages} séquentiels. La figure~\ref{fig:r2-s2-pipeline-stages}
présente la vue d'ensemble de ces étapes telles que définies dans \texttt{azure-pipelines.yml} :

\begin{enumerate}
  \item \textbf{Build} — compilation backend et frontend ;
  \item \textbf{CodeQuality} — analyse SonarQube ;
  \item \textbf{Docker} — construction des images ;
  \item \textbf{Push} — publication vers Nexus ;
  \item \textbf{Terraform} — provisionnement IaC ;
  \item \textbf{DeployK8s} — déploiement Kubernetes (Sprint~3, branche \texttt{main}) ;
  \item \textbf{DeployMonitoring} — Prometheus + Grafana via Ansible (Sprint~3).
\end{enumerate}

\begin{figure}[htbp]
    \centering
    \includegraphics[width=0.95\textwidth]{s2-pipeline-stages.png}
    \caption{Vue d'ensemble des stages du pipeline (\texttt{azure-pipelines.yml})}
    \label{fig:r2-s2-pipeline-stages}
\end{figure}

\subsubsection{Pool d'agents et connexion au dépôt}

Les tâches s'exécutent sur un \textbf{agent auto-hébergé} installé sur vm-gestion, regroupé
dans le pool \texttt{ticketmanager-pool}. La figure~\ref{fig:r2-s2-agent-pool} illustre la
configuration de ce pool d'agents. La figure~\ref{fig:r2-s2-repos} présente la structure du
dépôt Azure Repos, incluant les répertoires \texttt{backend}, \texttt{frontend},
\texttt{terraform} et \texttt{k8s}.

\begin{figure}[htbp]
    \centering
    \includegraphics[width=0.85\textwidth]{s2-agent-pool.png}
    \caption{Configuration du pool d'agents auto-hébergés}
    \label{fig:r2-s2-agent-pool}
\end{figure}

\begin{figure}[htbp]
    \centering
    \includegraphics[width=0.85\textwidth]{s2-azure-repos.png}
    \caption{Structure du dépôt Azure Repos (backend, frontend, terraform, k8s)}
    \label{fig:r2-s2-repos}
\end{figure}


% ---------------------------------------------------------------------------
\subsection{Réalisation}
% ---------------------------------------------------------------------------

\subsubsection{Compilation et tests}

\paragraph{Stage Build.}
Installation des dépendances npm (avec cache local sur vm-gestion), puis
\texttt{npm run build} pour le backend NestJS et le frontend Next.js.

\paragraph{Tests backend.}
Les tests unitaires et d'intégration s'exécutent avec Jest lorsque le stage Test est actif.
La figure~\ref{fig:r2-s2-tests-backend} présente le résultat de l'exécution des tests
automatisés du backend dans Azure DevOps.

\begin{figure}[htbp]
    \centering
    \includegraphics[width=0.85\textwidth]{s2-tests-backend.png}
    \caption{Exécution des tests automatisés du backend}
    \label{fig:r2-s2-tests-backend}
\end{figure}


\subsubsection{Analyse de la qualité du code (SonarQube)}

SonarQube analyse le code source backend (\texttt{backend/src}) et frontend
(\texttt{frontend}). La figure~\ref{fig:r2-s2-sonar-conn} présente la configuration de la
connexion entre SonarQube et Azure DevOps. La figure~\ref{fig:r2-s2-sonar-results} illustre
les résultats d'analyse du projet \texttt{ticket-manager}.

\begin{figure}[htbp]
    \centering
    \includegraphics[width=0.85\textwidth]{s2-sonar-connection.png}
    \caption{Connexion SonarQube — Azure DevOps}
    \label{fig:r2-s2-sonar-conn}
\end{figure}

\begin{figure}[htbp]
    \centering
    \includegraphics[width=0.85\textwidth]{s2-sonar-results.png}
    \caption{Résultats d'analyse SonarQube — projet TicketManager}
    \label{fig:r2-s2-sonar-results}
\end{figure}


\subsubsection{Construction des images Docker}

Deux images sont produites : \textbf{ticket-backend} (API NestJS, port 3001) et
\textbf{ticket-frontend} (Next.js, port 3000). La figure~\ref{fig:r2-s2-dockerfile-be}
présente le Dockerfile multi-étapes du backend. La figure~\ref{fig:r2-s2-dockerfile-fe}
détaille celui du frontend. La figure~\ref{fig:r2-s2-docker-build} montre la construction
réussie des images dans le stage Docker du pipeline.

\begin{figure}[htbp]
    \centering
    \includegraphics[width=0.85\textwidth]{s2-dockerfile-backend.png}
    \caption{Dockerfile du backend NestJS}
    \label{fig:r2-s2-dockerfile-be}
\end{figure}

\begin{figure}[htbp]
    \centering
    \includegraphics[width=0.85\textwidth]{s2-dockerfile-frontend.png}
    \caption{Dockerfile du frontend Next.js}
    \label{fig:r2-s2-dockerfile-fe}
\end{figure}

\begin{figure}[htbp]
    \centering
    \includegraphics[width=0.85\textwidth]{s2-docker-build.png}
    \caption{Construction des images dans le stage Docker}
    \label{fig:r2-s2-docker-build}
\end{figure}


\subsubsection{Scan de vulnérabilités (Trivy)}

Trivy analyse les images construites et signale les CVE par niveau de gravité. La
figure~\ref{fig:r2-s2-trivy-be} présente le rapport de scan pour l'image backend. La
figure~\ref{fig:r2-s2-trivy-fe} illustre celui de l'image frontend.

\begin{figure}[htbp]
    \centering
    \includegraphics[width=0.85\textwidth]{s2-trivy-backend.png}
    \caption{Scan Trivy — image backend}
    \label{fig:r2-s2-trivy-be}
\end{figure}

\begin{figure}[htbp]
    \centering
    \includegraphics[width=0.85\textwidth]{s2-trivy-frontend.png}
    \caption{Scan Trivy — image frontend}
    \label{fig:r2-s2-trivy-fe}
\end{figure}


\subsubsection{Publication vers Nexus Repository}

Nexus joue le rôle de \textbf{registre Docker privé} sur vm-gestion. La
figure~\ref{fig:r2-s2-nexus-repo} présente la création du dépôt Docker hosted. La
figure~\ref{fig:r2-s2-nexus-images} montre les images \texttt{ticket-backend} et
\texttt{ticket-frontend} stockées après exécution du pipeline.

\begin{figure}[htbp]
    \centering
    \includegraphics[width=0.85\textwidth]{s2-nexus-repo.png}
    \caption{Création du dépôt Docker hosted dans Nexus}
    \label{fig:r2-s2-nexus-repo}
\end{figure}

\begin{figure}[htbp]
    \centering
    \includegraphics[width=0.85\textwidth]{s2-nexus-images.png}
    \caption{Images \texttt{ticket-backend} et \texttt{ticket-frontend} dans Nexus}
    \label{fig:r2-s2-nexus-images}
\end{figure}


\subsubsection{Terraform automatisé dans le pipeline}

Le stage \textbf{Terraform IaC} exécute \texttt{terraform init}, \texttt{plan} et
\texttt{apply}. La figure~\ref{fig:r2-s2-tf-pipeline} présente l'exécution réussie de ce
stage dans Azure DevOps. La figure~\ref{fig:r2-s2-pipeline-run} illustre une exécution
complète du pipeline CI (stages Build à Push).

\begin{figure}[htbp]
    \centering
    \includegraphics[width=0.85\textwidth]{s2-terraform-pipeline.png}
    \caption{Exécution réussie du stage Terraform dans Azure DevOps}
    \label{fig:r2-s2-tf-pipeline}
\end{figure}

\begin{figure}[htbp]
    \centering
    \includegraphics[width=0.85\textwidth]{s2-pipeline-run-success.png}
    \caption{Exécution complète du pipeline CI (stages Build à Push)}
    \label{fig:r2-s2-pipeline-run}
\end{figure}


% ---------------------------------------------------------------------------
\subsection{Bilan du Sprint 2}
% ---------------------------------------------------------------------------

Le Sprint~2 met en place une chaîne d'intégration continue couvrant la compilation, la
qualité, la conteneurisation, la sécurité des images et la publication vers un registre
privé. Le déploiement effectif des pods applicatifs sur le cluster (stage
\textbf{DeployK8s}) et la supervision (stage \textbf{DeployMonitoring}) sont détaillés au
Sprint~3.
